% tipo di documento
\documentclass[12pt,a4paper,openright,oneside]{book}
\newenvironment{abstract}{}{}
\usepackage{abstract}
% codifica caratteri
\usepackage[utf8]{inputenc}
% encoding del testo
\usepackage[T1]{fontenc}
% dimensione dei margini
\usepackage[a4paper,top=2.5cm,bottom=2.5cm,left=3cm,right=3cm]{geometry}
% dimensione del font
\usepackage[fontsize=12pt]{scrextend}
% lingua del testo
\usepackage[english,italian]{babel}
% lingua per la bibliografia
\usepackage[fixlanguage]{babelbib}
% package per generare testo fittizio. Potrebbe essere
% utile nel controllare quanto un capitolo potrebbe essere
% grande e quindi quanto occupa nella pagina
\usepackage{lipsum}
% per ruotare le immagini
\usepackage{rotating}
% per modificare l'header delle pagine 
\usepackage{fancyhdr}
% per allineare in modo giustificato
\usepackage{ragged2e}
\justifying
% uso delle immagini
\usepackage{graphicx}
\usepackage{float}
% uso dei colori
\usepackage[dvipsnames, table]{xcolor}         
% uso dei listing per il codice
\usepackage{listings}       
% per inserire gli hyperlinks tra i vari elementi del testo 
\usepackage[colorlinks=true, allcolors=black]{hyperref}    \usepackage{verbatim}
% diversi tipi di sottolineature
\usepackage[normalem]{ulem}
% package e comando per creare pagine vuote
\usepackage{afterpage}
\newcommand\blankpage{%
    \null
    \thispagestyle{empty}%
    \addtocounter{page}{-1}%
    \newpage
}
    
% package per creare comandi personalizzati
\usepackage{xpatch}
% package helper per le liste puntate
\usepackage{enumitem}
% package per l'utilizzo dei colori
% package per l'highlighting del codice
\usepackage{minted}
% package per gestire le caption
\usepackage{caption}
\usepackage{subcaption}
% per gestire tabelle su più pagine
\usepackage{longtable}
% per combinare le righe di una tabella
\usepackage{multirow}
% per creare i tree di directory
\usepackage{dirtree}
\usepackage{adjustbox} % Aggiungi questo nel preambolo
\usepackage{graphicx}


% per le icone a fianco dei titoli di sezione
\usepackage{etoolbox}
\newcommand{\icon}[1]{\includegraphics[height=12pt]{#1}}
\robustify{\icon}

% -----------------------------------------------------------------

% Modifica lo stile dell'header
\pagestyle{fancy}
\fancyhf{}
\lhead{\rightmark}
\rhead{\textbf{\thepage}}
\fancyfoot{}
\setlength{\headheight}{15pt}

% Rimuove il numero di pagina all'inizio dei capitoli
\fancypagestyle{plain}{
  \fancyfoot{}
  \fancyhead{}
  \renewcommand{\headrulewidth}{0pt}
}

% comandi per cambiare temporaneamente la lingua
% abstract in inglese, al fine di cambiarne il titolo
\xpretocmd{\abstract}{\selectlanguage{english}}{}{} 
\xapptocmd{\endabstract}{\selectlanguage{italian}}{}{}

% formattazione e highlight del codice
\usemintedstyle{manni}
\setminted[typescript]{
    framesep=2mm,
    baselinestretch=1.2,
    fontsize=\ttfamily\footnotesize,
    linenos
}

% rimozione del prefix per le tabelle
\captionsetup[table]{labelformat=empty}

% environment per impostare il codice in piu' pagine
\newenvironment{code}{\captionsetup{type=listing}}{}

% -----------------------------------------------------------------
\begin{document}

\begin{titlepage}

\begin{center}
    \textbf{\huge{Università degli Studi di Torino}}
    \vspace{2mm}
    \\ \LARGE{Dipartimento di Informatica}
    \vspace{5mm}
    \\ \includegraphics[keepaspectratio=true,width=5cm]{images/unito_logo.pdf}
    \vspace{5mm}
\end{center}

\begin{center}
    \LARGE{Tesi di Laurea Magistrale in Informatica} 
\end{center}

\vspace{15mm}
\begin{center}
    \textbf{\huge{ Accessibilità dei servizi primari}}
\end{center}
\vspace{30mm}

\begin{minipage}[t]{0.47\textwidth}
	{\large{Relatore:}{\normalsize\vspace{3mm}
	\bf\\ \large{Rossano Schifanella} \normalsize\vspace{3mm}\bf}}
\end{minipage}
\hfill
\begin{minipage}[t]{0.47\textwidth}\raggedleft
	{\large{Candidato:}{\normalsize\vspace{3mm} \bf\\ \large{Martina Celli}}}
\end{minipage}

\vspace{40 mm}
\hrulefill
\\ \centering{\large{ANNO ACCADEMICO 2026/2027}}

\end{titlepage}

\thispagestyle{plain}

{
\raggedleft
\textit{Chi è deserto non vuole che qualcosa fiorisca in te}
\\ - 883 Ci sono anche io

}
\include{contents/1.Abstract}
\thispagestyle{plain}
\vspace*{\fill}
\noindent\textit{Dichiaro di essere responsabile del contenuto dell’elaborato che presento al fine del conseguimento del titolo, di non avere plagiato in tutto o in parte il lavoro prodotto da altri e di aver citato le fonti originali in modo congruente alle normative vigenti in materia di plagio e di diritto d’autore. Sono inoltre consapevole che nel caso la mia dichiarazione risultasse mendace, potrei incorrere nelle sanzioni previste dalla legge e la mia ammissione alla prova finale potrebbe essere negata.}
\vspace*{\fill}

\tableofcontents




\include{contents/SampleChapter}


\bibliography{Bibliography}
\bibliographystyle{plain}

\include{contents/Thanks}

\end{document}
% -----------------------------------------------------------------
